\section{Discussion: Open Networking Questions}
\label{sec:discussion}

Treating self-regulating cars as a deployment problem turns their hardest questions into systems and networking questions.

\noindent\textbf{How little communication is enough, and can it be private?}
Section~\ref{sec:vision} fixes semantics, not sizes. We need principled positions on rate, precision, aggregation radius, freshness, and trust: how stale a signal may be before it misleads, how coarse before it stops helping, and how a receiver should weight summaries of unknown provenance. Privacy sharpens the question. Our contract keeps the consumed unit anonymous, but a naive realization still has each sender emit its own position, as a Basic Safety Message already does; the open problem is to produce the aggregate without any raw per-vehicle trajectory leaving the vehicle, through in-network or gossip-style aggregation, or randomized-response and differentially private counts that stay useful to a flow regulator yet carry no linkable identity.

\noindent\textbf{What belongs in the cloud, and how late may it be?}
The runtime loop is fully local and offline-capable by construction; what remains is training, periodic model updates, and fleet analytics over decimated uplinks of a kilobyte per vehicle-second or less. Delay tolerance is thus a design property rather than a hope: an advisory never waits on backhaul. The open question is how much slowly arriving global context, incidents or corridor demand, is worth folding back into the local policy, and on what cadence.

\noindent\textbf{How robust is the system to partial deployment?}
Evaluation should sweep penetration, compliance, message loss, sensing noise, and demand bursts across a topology ladder from rings to merges to inverted trees, should directly compare bottleneck seeding against uniform adoption at equal fleet size, and should report throughput, jam fraction, and delay alongside hard braking, collisions, and follower delay, so that a network gain is never bought by making human traffic less safe.

\noindent\textbf{How should safety be verified?}
A local policy can optimize a network objective and still be unacceptable if it surprises human drivers. The safety layer should be specified, tested, and audited separately from the learned controller, with explicit checks for stale state, follower disruption, low-speed obstruction, and unsafe merges. This separation can become a systems requirement: the executed-action diagnostics of Section~\ref{sec:architecture} report how often a suggestion was masked, blocked, or overridden, giving continuous evidence that a network objective is being achieved safely.

\noindent\textbf{Who benefits, and who decides?}
A voluntary advisory system can shift delay across drivers, lanes, or vehicle classes. A deployable system therefore needs measurement and incentives that make the network benefit visible and its distribution accountable, especially in the fleet-seeded model, where one operator acts and every driver rides the smoother flow.

As vehicles become networked computers, congestion control need not remain a roadside service: the controller fits in coarse bins, the radio is already deployed, the compute costs a few hundred dollars, and a single commercial fleet is enough presence at a real bottleneck. The open work is to give each vehicle just enough local traffic state to help the network while it acts independently, safely, privately, and without a cloud coordinator; we think this is a productive problem space for the networking community.
